% ============================================================================
%  引用与链接：preamble/references-setup.tex
%  ★ 必须最后加载：hyperref 之后再加载 cleveref，否则交叉引用名称会错乱。
% ============================================================================

\definecolor{linkcolor}{RGB}{0, 102, 204}
\definecolor{citecolor}{RGB}{0, 128, 0}
\definecolor{urlcolor}{RGB}{85, 26, 139}

\usepackage{hyperref}
\hypersetup{
  colorlinks     = true,
  linkcolor      = linkcolor,
  citecolor      = citecolor,
  urlcolor       = urlcolor,
  pdfstartview   = {Fit},
  bookmarksnumbered = true,
  bookmarksopen  = true,
}

% pdftitle / pdfauthor 依赖 config/metadata.tex 中的宏，故延后到开文时才展开
\AtBeginDocument{%
  \hypersetup{pdftitle={\papertitle}, pdfauthor={\authorname}}%
}

\usepackage[nameinlink]{cleveref}

% --- cleveref 中文名称 ---
\crefname{section}{节}{节}
\Crefname{section}{节}{节}
\crefname{figure}{图}{图}
\Crefname{figure}{图}{图}
\crefname{table}{表}{表}
\Crefname{table}{表}{表}
\crefname{equation}{式}{式}
\Crefname{equation}{式}{式}
\crefname{algorithm}{算法}{算法}
\Crefname{algorithm}{算法}{算法}
\crefname{appendix}{附录}{附录}
\Crefname{appendix}{附录}{附录}
\crefname{theorem}{定理}{定理}
\crefname{lemma}{引理}{引理}
\crefname{definition}{定义}{定义}

% --- 语义化引用快捷命令（推荐优先于裸 \ref）---
% 注意：不要取 \reffig/\reftab/\refalg 等名字，它们可能与 cleveref 自动生成的
%       \<float>ref 系列冲突（实测 \refalg 会被 cleveref 占用）。统一用 \ref* 前缀。
\newcommand{\reffig}[1]{图~\ref{#1}}
\newcommand{\reftab}[1]{表~\ref{#1}}
\newcommand{\refsec}[1]{第~\ref{#1}~节}
\newcommand{\refeqn}[1]{式~(\ref{#1})}
\newcommand{\refalg}[1]{算法~\ref{#1}}
\newcommand{\refapp}[1]{附录~\ref{#1}}

% 说明：elsarticle 采用数字制引文，直接使用 \cite{} 即可（多文献自动压缩为 [1,3--5]）。
